\documentclass{article}

\expandafter\let\csname equation*\endcsname\relax
\expandafter\let\csname endequation*\endcsname\relax

\usepackage[dvipsnames,table]{xcolor}  % Coloured text etc.
\usepackage{natbib}
\setcitestyle{square,numbers}
\usepackage{hyperref}
\newcommand\myshade{85}
\colorlet{mylinkcolor}{violet}
\colorlet{mycitecolor}{YellowOrange}
\colorlet{myurlcolor}{Aquamarine}
\hypersetup{
    colorlinks=true,
    linkcolor  = mylinkcolor!\myshade!black,
    citecolor  = mycitecolor!\myshade!black,
    urlcolor   = myurlcolor!\myshade!black,
}

% \usepackage[pdftex]{graphicx} 
\usepackage[shortlabels]{enumitem}
\usepackage{caption}

\captionsetup{
    format=plain,
    justification=RaggedRight,
    singlelinecheck=true
}

\usepackage{graphicx}
\usepackage{xargs}
\usepackage{subcaption}
\usepackage{multirow}
\usepackage{amsmath}
\usepackage{amssymb}
\usepackage{mathtools}
\usepackage{float}
\usepackage{nicefrac}
\usepackage{pifont}
\usepackage{booktabs}
\usepackage[a4paper, total={6in, 8in}]{geometry}
\newtheorem{lemma}{Lemma}
\newtheorem{corollary}{Corollary}
\newtheorem{remark}{Remark}
\newtheorem{proposition}{Proposition}

\newcommand{\x}{\mathbf{x}}
\renewcommand{\u}{\mathbf{u}}
\renewcommand{\v}{\mathbf{v}}
\renewcommand{\b}{\mathbf{b}}
\newcommand{\p}{\mathbf{p}}
\newcommand{\q}{\mathbf{q}}
\newcommand{\A}{\mathbf{A}}
\newcommand{\V}{\mathbf{V}}
\newcommand{\w}{\mathbf{w}}
\newcommand{\s}{\mathbf{s}}
\newcommand{\y}{\mathbf{y}}
\renewcommand{\H}{\mathbf{H}}
\newcommand{\B}{\mathbf{B}}
\newcommand{\Hnaut}{\mathbf{H}_0}
\newcommand{\Bnaut}{\mathbf{B}_0}
\newcommand{\U}{\mathbf{U}}
\newcommand{\prox}{\operatorname{prox}}
\newcommand{\argmin}{\operatorname{argmin}}
\newcommand{\argmax}{\operatorname{argmax}}

\graphicspath{{fig/}}

\begin{document}
\title{Scalable Proximal QuasiNewton Methods for Collision Resolutions of Dense Suspensions}
\author{{Nic Rummel}, {Eduardo Corona}}
\date{\today}
\maketitle

\abstract{Collision-resolution is key subalgorithm necessary in simulating suspensions of rigid particles in Newtonian Stokes flow. Our approach formulates the collision resolution problem as a linear complementarity problem with geometric ‘non-overlapping’ constraints imposed at each time-step which is then reformulated as a constrained quadratic program (QP). We develop a novel proximal QuasiNewton method which takes advantage of the structure of the QP to provide convergence guarantees without the need for a linesearch with no extra matrix-vector products. The algorithm significantly outperforms the current state of the art gradient based methods.}

\section{Background}
At every time iteration of the time evolution of system of interest, $N$ particle that lie in a threshold are identified as possibly colliding within the time stepping window. This motivates solving the following a linear complimentary problem
\[\A\x -\b \geq \mathbf{0} \perp \x \geq \mathbf{0}\]
where $\A \in \mathbb{S}_{++}^{n \times n}, \b \in \mathbb{R}^n$. While $\A$ is the mathematically symmetric and positive definitedue to numerical error it can lose symmetry. From our testing, because of the spectrally accurate evaluations of the the boundary integral solvers the lose of symmetry is negligible. Depending on how the problem is set up $n=N$ or $n=3N$. In the later case, each possible contact is given three degrees of freedom to allow for non-tangential forces. The values of $\x$ either satisfy a linear system or are element-wise 0. Physically $\x$ corresponds to forces and torques acting on individual particle from a collision. If the particles collide, then the forces are nonzero, but if the particles do not collide then the collisional forces are zero.

% \textcolor{green}{Eduardo: The size of $A$ is incorrect. $N$ is either the number of contacts (frictionless) or $3$ times the number of contacts. $A$ is a product of matrices which are bigger in size, but that needs to be specified elsewhere.}

% \textcolor{blue}{Nic: Is that better? In my preliminary testing I believe $n=N$, but my reasoning for why that is could be wrong.}
% \textcolor{red}{Stephen: is $\A$ symmetric? PSD? PD?}

% \textcolor{blue}{Nic: Mathematically is PD $\A \succ 0$, but computational errors in the evaluation of the boundary integrals could cause $\A$ to loose symmetry. Thus, even if we wanted to use a non-neg CG solver we would need to use a Krylov method that does not rely on symmetry.} 

For small sized systems the LCP can be solved directly with pivoting  methods. For medium sized systems or interior point methods and minimum map Newton are competitive choices \citep{NiebeErleben2015}. When the system becomes large enough these methods become computationally infeasible \cite{yanScalableComputationalPlatform2020}. Another approach is to notice that satisfying the LCP is the KKT conditions of the following constrained quadratic program. 
\begin{equation} \label{eq:qp}
    \x^* = \argmin_{\x \geq 0 }  f(\x) =:\bigl\langle \x, \tfrac{1}{2}\A\x +\b\bigr\rangle
\end{equation}

Due to the size of systems of interests, it is not computationally efficient to store the matrix $\A$.  Thus, only a matrix-vector product is available. Each mat-vec involves solving from boundary integrals on each particle. These integrals solve the system of partial differential equations to obtain forces of the fluid flow acting on each particle. 

This motivate finding optimization problems that solve Equation \eqref{eq:qp} with as few applications of the matrix $\A$ as possible. For large systems, it is standard to apply projected gradient descent (PGD). The choice of search direction comes from approximating the cost function with the following quadratic model:
\[ \x_{k+1} = \underset{\x \geq 0}{\argmin} \langle \A \x_k + \b, \x-\x_k \rangle + \tfrac{L}{2}\langle \x-\x_k , \x-\x_k \rangle\]
where $L = \|\A\|_2$. Obtaining $L$ requires a power method that utilizes as many evaluations of $\A\cdot$ as the PGD algorithm itself. Implicitly, the Hessian is being approximated by a scaled identity. From standard convergence theory, we know that the rate of convergence in the unconstrained case is $\mathcal{O}(\kappa(\A)N /\epsilon)$ where $\kappa(\A)$ is the condition number of the matrix $\A$ and $\epsilon$ is the desired tolerance of the approximation \citep{Atkinson1989}. In the constrained case, we can solve this problem in closed form. This leverages that the proximal mapping is separable because the approximation to the Hessian is diagonal. The stepping rule applies orthogonal projection onto the positive orthant of the standard gradient descent step
\[ \x_{k+1} = \max\bigl(\mathbf{0}, \x_k - \kappa (\A \x_k + \b)\bigr)\]
for some step size $ \kappa \in (0,L]$.

The convergence can be improved by leveraging the constant Hessian of the QP, thus using curvature information should bring significant improvements on the rate of convergence. Unfortunately, because $\A$ is not explicitly available Newton like algorithms become infeasible. Specifically, application of $\A$ would require the use a Krylov subspace method such as conjugate gradient, and this makes the cost of each iteration prohibitive. Instead, we suggest approximating the pertinent curvature information with a QuasiNewton method. Specifically, we solve the following  subproblem at the current iterate 
\begin{equation} \label{eq:qn-quadMod}
    \x_{k+1} = \underset{\x \geq 0}{\argmin} \langle \A \x_k + \b, \x-\x_k \rangle + \tfrac{1}{2}\langle \x-\x_k , \B_k(\x-\x_k )\rangle %= \prox^{\B_k}_{\cdot \geq \mathbf{0}}(\x_k + \H_k )
\end{equation}
where $\B_k$ is a current approximation of $\A$ given from QuasiNewton updates. Because $\B_k$ is obtained from formula that only require information about the gradient at the current and past iterates there is no increase in the number of applications of $\A$ per iteration.% In particular, in the unconstrained case this increases the cost of the each iteration by $\mathcal{O}(nr)$ where $r$ is the memory of the quasi-Newton approximation. 

Due to the non-negativity constraint and constraints this subproblem cannot be solved in closed form for a general $\B_k$. Schmidt et al. suggested solving the subproblem with variant of PGD \citep{SchmidtBergFriedlanderEtAl2009ProcTwelfthIntConfArtifIntellStat}, and Kim et al. suggested leveraging a linesearch on the original cost function \eqref{eq:qp} to guarantee the armijo condition and using binding sets to modify the $\B_k$ in such a way that stepping into the infeasible region happens infrequently \citep{KimSraDhillon2010SIAMJSciComput}. Becker et al. takes a different route and leverages the structure of $\B_k = \B_0 + \U\U^\top - \V\V^\top$ where $\U \in \mathbb{R}^{n\times r_1}$ and $\V \in \mathbb{R}^{n\times r_2}$. This diagonal and low rank structure gives rise to an alternative form of the subproblem  
\begin{equation} \label{eq:prox-step}
    \begin{split}
        \y &=: \x_k + \H_k (\A\x_k + \b)\\
        \mathbf{C} &=: \B_0 + \U\U^\top\\
         \x_{k+1} &= \operatorname{max}\left(\mathbf{0},\y + \begin{bmatrix} 
            - \B_0^{-1} U & \mathbf{C}^{-1} \V 
        \end{bmatrix} \boldsymbol{\alpha}^*\right) 
\end{split}
\end{equation}
where $\boldsymbol{\alpha}^* \in \mathbb{R}^{2r}$ is the unique root of piecewise linear function
\[
\mathcal{L}\left(\boldsymbol{\alpha}\right) = \begin{bmatrix} 
        \U^{\top}&0 \\ 0&\V^{\top} 
    \end{bmatrix} \begin{bmatrix} 
       \y+[\mathbf{0}_{r\times r}, \mathbf{C}^{-1} \V] \boldsymbol{\alpha}- \operatorname{max}\left(\mathbf{0},\y + \begin{bmatrix} 
            - \B_0^{-1} U & \mathbf{C}^{-1} \V 
        \end{bmatrix} \boldsymbol{\alpha}\right) \\
       \y -  \operatorname{max}\left(\mathbf{0},\y + \begin{bmatrix} 
            - \B_0^{-1} U & \mathbf{C}^{-1} \V 
        \end{bmatrix} \boldsymbol{\alpha}\right)  
    \end{bmatrix} + \boldsymbol{\alpha}.
\]
This is a direct application the general convex analysis result given in Corollary 3.6 of \citep{BeckerFadiliOchs2018arXiv180108691}. Becker et al. developed this approach for general optimization problems where the objective function could be split into a smooth function plus a simple but not smooth function. The leveraged this result to implement fast zero memory quasi-newton methods. 

Specifically, when the Quasi-Newton update rule is chosen to be the SR1 update with zero memory finding the root of $\mathcal{L}(\alpha)$ is a one dimensional problem and can be solved in $\mathcal{O}(n\log(n))$ time by locating the interval where the root occurs and solving for the root in closed form, see Appendix \ref{sec:0SR1}. In the general case, the root can be found using a semi-smooth newton solver, see Appendix \ref{sec:ss-newton}. The benefit of this approach of is that root finding problem is significantly smaller than the original problem. Thus the computational burden of the overall algorithm lies in the number of iterations.

\section{Step Size}
For the unconstrained case, once a search direction $\p_k$ is chosen an optimal step size can be computed with the cost of two mat-vec's: 
% \[\kappa^* = \argmin_\kappa \bigl\langle \x_k - \kappa\p_k, \tfrac{1}{2}\A(\x_k - \kappa\p_k) +\b\bigr\rangle = -\frac{\langle \nabla f_k, \p_k\rangle }{\langle \p_k,\A\p_k\rangle}\]
% So for the price of one extra mat-vec the optimal step size can be computed explicitly. 
\begin{align*}
    \kappa^* &= \argmin_\kappa \langle \x + \kappa\p, \tfrac{1}{2}\A(\x + \kappa\p) +\b\rangle \\
    0 &\stackrel{set}{=} \frac{d}{d\kappa } \bigl(\langle \x + \kappa\p, \tfrac{1}{2}\A(\x + \kappa\p) +\b\rangle\bigr)\\
    % &=\frac{d}{d\kappa } \bigl(\langle \x + \kappa\p, \tfrac{1}{2}\A(\x + \kappa\p) \rangle + \langle\x + \kappa\p,\b\rangle\bigr)\\
    % &=\frac{d}{d\kappa } \bigl(\langle \x, \tfrac{1}{2}\A(\x + \kappa\p) \rangle + \tfrac{1}{2} \langle \kappa\p,\A\x\rangle  + \tfrac{1}{2} \langle \kappa \p, \A \kappa\p\rangle + \langle\x + \kappa\p,\b\rangle\bigr)\\
    % &= \tfrac{1}{2} \langle \x, \A \p\rangle +\tfrac{1}{2} \langle \p, \A \x\rangle +  \kappa \langle \p, \A \p\rangle + \langle \p, \b\rangle\\
    \kappa^* = h(\p) &=: -\frac{\tfrac{1}{2} (\langle \x, \A \p\rangle + \langle \p, \A \x\rangle) + \langle \p, \b\rangle}{\langle \p, \A \p\rangle}
\end{align*}

In the constrained case, the following iteration could be used 
\begin{align*}
    \p_k &= \H_k \nabla f_k &\text{(Quasi-Newton Step Direction)}\\
    \kappa_k &= h(\p_k) &\text{(Unconstrained Optimal Step Size)}\\
    \bar{\x}_k &= \prox_{\x\geq 0}^{\B_k}(\x_k +\kappa_k \p_k) &\text{(Apply Proximal Operator)}\\
    \q_k &= \bar{\x}_k - \x_k &\text{(Feasible Step Direction)}\\
    \eta_k &= \min\left(1,  h(\q_k)\right) &\text{(\emph{Semi}-Optimal Feasible Step Size)} \\
    \x_{k+1} &= \x_k + \eta^* \q_k &\text{(Next Iterate)} 
\end{align*}
More testing and analytic compution is require to demonstrate whether the additional mat-vecs per iteration are necessary. 

In the literature it has been shown for quasi-Newton methods (without limited memory) that $\alpha^* \rightarrow 1$. As we are in the limited memory regime it is less clear as wether this will be the case. 
\section{Stopping Criterion}
A nature merit function is 
\begin{equation}\label{eq:stop-crit}
    \mathrm{e}_{\text{kkt}}=\tfrac{1}{2}\|\phi\|^2_2
\end{equation}
where $\phi = \bigl(\min(x_i, \nabla f_i)\bigr)_{i=1}^n$. This corresponds to how accurately the LCP is satisfied under the Euclidean norm, or equivalently the KKT conditions of the quadratic program. We refer to absolute error as $\mathrm{e}_{\text{kkt}}$ and relative error as
\[\mathrm{e}_{\text{rel kkt}} = \frac{\left|\mathrm{e}_{\text{kkt}}^{(k)} - \mathrm{e}_{\text{kkt}}^{(k-1)}\right|}  {\mathrm{e}_{\text{kkt}}^{(k)}}.\]


\section{Hessian Approximation}

At each iteration the distance between the two most recent iterates and gradients are specified as 
\[\s = \x_k - \x_{k-1} \;\mathrm{ and } \;\y = \nabla f_k - \nabla f_{k-1} \;\mathrm{respectively.}\]
The goal is at each iteration is improve the approximation to the inverse Hessian $\H \approx \nabla^2 f_k ^{-1}$ that satisfies the secant condition $\H \y = \s$.

In the unconstrained case, we only need to apply the matrix vector product $\H_k( \cdot)$. There are different approaches, but in the limited memory case it is standard to save the vectors $\{\s_j, \y_j\}_{j=k-r}^{k}$. This allows for a relatively fast matrix vector product computation of $O(n r)$, but for the proximal operator we need to represent the Hessian approximation as a diagonal plus a rank $r$ update $\B = \B_0 + \U\U^{\top} - \V\V^{\top}$. 
% If we choose to use the SR1 update we have that 
% \begin{align*}
%     \mathbf{S} &= [\s_{k-r} \; \cdots\;\s_k],\;
%     \mathbf{Y} = [\y_{k-r}\;\cdots \;\y_k]\\
%     (\mathbf{A})_{ij} &= 
%         \langle \s_{k-r+i-1}, \y_{k-r+j-1}\rangle\\
%     \B_k &= \B_0 +\left(\mathbf{Y}-\B_0 \mathbf{S}\right)\left(\mathbf{A}-\mathbf{S}^{\top} \B_0 \mathbf{S}\right)^{-1}\left(\mathbf{Y}-\B_0 \mathbf{S}\right)^{\top},\\
% \end{align*}

\subsection{SR1}
The symmetric rank 1 (SR1) update is 
\[\gamma \in (0,1], \; \tau_{BB_2} = \frac{\langle\s, \y\rangle}{\|\y\|^2} > 0, \; h_0 = \gamma \tau_{BB_2},\; \Hnaut = h_0 I, \]
\[\u = \frac{\s - \Hnaut \y}{ \sqrt{\sigma \langle \s - \Hnaut \y, \y\rangle}},\;\H = \Hnaut + \sigma \u \u^{\top}.\; \]
where $\sigma = \operatorname{sign}[\langle \s - \Hnaut \y, \y\rangle] \in \{-1,1\}$. This allows for the update to either be the addition or subtraction of a rank one matrix.Nocedal and Wright (page 145) specify that SR1 is compelling for the following reasons
\begin{itemize}
    \item There are simple fast ways to check for numerical break downs of the method compared to BFGS (which requires damping in the constrained case). 
    \item Because we are minimizing a quadratic, the Hessian of the objective is constant, and the Hessian approximations generated by the SR1 are often better than the BFGS approximations.
    \item For a constrained problem, imposing the curvature condition  $y_k^{\top} s_k>0$ is often impossible, and thus BFGS updating is not recommended. Indefinite Hessian approximations are desirable insofar as they reflect indefiniteness in the true Hessian.
\end{itemize}
\subsection{BFGS}
$\Hnaut$ is chosen in the same way as for the SR1 update, and thus $\Bnaut$ is specified as well. 

This is a rank-2 update at every step: 
\[\B_{k+1} = \B_k + \u_k\u^{\top}_k - \v_k\v^{\top}_k\]
\[\u_k = \frac{\y_k }{\sqrt{ \langle \y_k,\s_k\rangle}}, \;\v_k = \frac{\B_k \s_k }{\sqrt{ \langle \s_k, \B_k \s_k \rangle}}\]
There are some 
\begin{itemize}
    \item BFGS guarantees that the Hessian approximation stays positive definite
    \item No need for a trust region subproblem at each iteration
\end{itemize}






\subsection*{Preliminary Results}
The initial experiments was to save out the LCP's created from running `RBS\_Mobility.m' with on four different scenarios 
\begin{itemize}
    \item 2x2x2 lattice of bimetallic particles.
    \item 3x3x3 lattice of bimetallic particles.
    \item 4x4x4 lattice of amphiphilic particles.
    \item 5x5x5 lattice of amphiphilic particles.
\end{itemize}. 
All solvers were given the tolerance of 1E-12 for absolute error on the KKT condition and 1E-12 of relative error of the KKT condition, but the error information for each iteration was stored for all solvers with a callback function. This allows us to plot the results at different error tolerances. The metric of particular interest is the number of matrix vector products as each corresponds to a boundary integral solve. 

All the saved LCP's were separated into the corresponding size: $n \in [10,50], n\in [50,100], n \in [100,150]$ and $n \in [150,200]$. QuasiNewton methods were ran with a limited memory context window of 20. No line searches were used for any algorithms except for L\_BFGS\_B. The tolerance of 1E-6 for both absoluted and 
% \subsection{Non-Symmetric System}

% Here we are solving 
% \[\min_{\x\geq 0} \tfrac{1}{2}\|\A\x+b\|_2^2\]
% This has the gradient of 
% \[\nabla f = \A^\top(\A\x +\b)\]
% and KKT conditions of 
% \[0 = \A^\top(\A\x +\b) \perp \x = 0 \]
% When $\A$ is full rank (which it is) then this is equivalent to 
% \[0 = \A^\top(\A\x +\b) \perp \x = 0 \]
% So we can solve this non-negative least squares problem.
% \begin{figure}[H]
%     \centering
%     \includegraphics[width=\textwidth]{n_10_50_boxPlot.pdf}
%     \caption{\protect \input{fig/n_10_50_caption.tex}}
%     \label{fig:bi-easy}
% \end{figure}

% \begin{table}[ht]
%     \resizebox{\textwidth}{!}{
%     \input{fig/n_10_50_table.tex}
%     }
%     \caption{Table for $n\in[10, 50]$}
% \end{table}

% \begin{figure}[H]
%     \centering
%     \includegraphics[width=\textwidth]{n_50_100_boxPlot.pdf}
%     \caption{\protect \input{fig/n_50_100_caption.tex}}
%     \label{fig:bi-med}
% \end{figure}

% \begin{table}[ht]
%     \resizebox{\textwidth}{!}{
%     \input{fig/n_50_100_table.tex}
%     }
%     \caption{Table for $n\in[50, 100]$}
% \end{table}

% \begin{figure}[H]
%     \centering
%     \includegraphics[width=\textwidth]{n_100_150_boxPlot.pdf}
%     \caption{\protect \input{fig/n_100_150_caption.tex}}
%     \label{fig:ammphi-easy}
% \end{figure}

% \begin{table}[ht]
%     \resizebox{\textwidth}{!}{
%     \input{fig/n_100_150_table.tex}
%     }
%     \caption{Table for $n\in[100, 150]$}
% \end{table}

% \begin{figure}[H]
%     \centering
%     \includegraphics[width=\textwidth]{n_150_200_boxPlot.pdf}
%     \caption{\protect \input{fig/n_150_200_caption.tex}}
%     \label{fig:ammphi-med}
% \end{figure} 

% \begin{table}[ht]
%     \resizebox{\textwidth}{!}{
%     \input{fig/n_150_200_table.tex}
%     }
%     \caption{Table for $n\in[150, 200]$}
% \end{table}

% \pagebreak

\subsection{Symmetric System}

Here we are solving 
\[\min_{\x\geq 0} \langle\x, \tfrac{1}{2}\hat{\A}\x+b\rangle\]
where $\hat{\A} = \tfrac{1}{2}(\A+\A^\top)$. This has the gradient of 
\[\nabla f = \hat{\A}\x +\b\]
and KKT conditions of 
\[0 = \hat{\A}\x +\b \perp \x = 0 \]
This does not solve the same problem as the LCP, but it is close, and mathematically $\A$ should be symmetric.
\begin{figure}[H]
    \centering
    \includegraphics[width=\textwidth]{optBwd_n_10_50_boxPlot.pdf}
    \caption{\protect \input{fig/optBwd_n_10_50_caption.tex}}
    \label{fig:bi-easy}
\end{figure}

\begin{table}[ht]
    \resizebox{\textwidth}{!}{
    \input{fig/optBwd_n_10_50_table.tex}
    }
    \caption{Table for $n\in[10, 50]$}
\end{table}

\begin{figure}[H]
    \centering
    \includegraphics[width=\textwidth]{optBwd_n_50_100_boxPlot.pdf}
    \caption{\protect \input{fig/optBwd_n_50_100_caption.tex}}
    \label{fig:bi-med}
\end{figure}

\begin{table}[ht]
    \resizebox{\textwidth}{!}{
    \input{fig/optBwd_n_50_100_table.tex}
    }
    \caption{Table for $n\in[50, 100]$}
\end{table}

\begin{figure}[H]
    \centering
    \includegraphics[width=\textwidth]{optBwd_n_100_150_boxPlot.pdf}
    \caption{\protect \input{fig/optBwd_n_100_150_caption.tex}}
    \label{fig:ammphi-easy}
\end{figure}

\begin{table}[ht]
    \resizebox{\textwidth}{!}{
    \input{fig/optBwd_n_100_150_table.tex}
    }
    \caption{Table for $n\in[100, 150]$}
\end{table}

\begin{figure}[H]
    \centering
    \includegraphics[width=\textwidth]{optBwd_n_150_200_boxPlot.pdf}
    \caption{\protect \input{fig/optBwd_n_150_200_caption.tex}}
    \label{fig:ammphi-med}
\end{figure} 

\begin{table}[ht]
    \resizebox{\textwidth}{!}{
    \input{fig/optBwd_n_150_200_table.tex}
    }
    \caption{Table for $n\in[150, 200]$}
\end{table}
\pagebreak

\bibliographystyle{abbrv}
\bibliography{CollisionResolution}
\appendix 
\section{Appendix}
\input{appendix.tex}
\end{document}